\documentclass[12pt]{article}
\usepackage[utf8]{inputenc}
\usepackage[T1]{fontenc}
\usepackage{hyperref}
\title{Ravula Language Resources for K-12 Curriculum Development}
\author{Compiled by Claude Sonnet 5.5 (claude-sonnet-5-5) for Andr\'as Kornai}
\date{}
\begin{document}
\maketitle

\section{Introduction}
Ravula (ISO 639-3: yea; Glottocode ravu1237) is a Dravidian language of southern India. Its speakers are called Yerava in Karnataka and Adiyan in Kerala. The 2011 census recorded about 26,563 speakers, and Joshua Project gives about 27,000. Most live in Kodagu district, Karnataka (about 25,000), with the rest in Wayanad district, Kerala. It is placed among the Malayalamoid languages. It is written in the Kannada script in Karnataka and the Malayalam script in Kerala. Sources differ on its endangerment status; Wikitongues lists it as EGIDS 5.

\section{Primary Educational Resources}
No primers or textbooks were found in search results.

\section{Linguistic and Reference Materials}
\begin{itemize}
\item Wikipedia, Ravula language: \url{https://en.wikipedia.org/wiki/Ravula_language}
\item Omniglot, Ravula (writing and background): \url{https://omniglot.com/writing/ravula.htm}
\item Wikitongues, Ravula (yea): \url{https://wikitongues.org/languages/yea/}
\end{itemize}

\section{Religious and Community Resources}
\begin{itemize}
\item Joshua Project, Yerava in India. It reports that the New Testament and the JESUS Film exist in Ravula: \url{https://joshuaproject.net/index.php/people_groups/18343/IN}
\end{itemize}

\section{Digital and Audio Resources}
No dedicated audio or software resources were identified apart from the Joshua Project listing above and the Wikitongues language page.

\section{Curriculum Development Context}
Ravula has a sizeable community split across two states and two scripts (Kannada and Malayalam), so any curriculum must settle on a script. The dominant school languages there are Kannada, Malayalam and English. No evidence of Ravula-medium schooling or published teaching materials was found. A first step would be to locate the New Testament text and any community documentation as seed corpora.
\end{document}